\chapter[demand_research]{Activités de recherche}
 [
  Ce chapitre résume mes activités de recherche passées, présentes et futures. \\
  Tout le texte de ce document est un texte de remplacement : remplacez-le par le vôtre. \\
  Pour des raisons de place, les travaux tiers ne sont pas listés en fin de document. Vous trouverez donc çà et là des références à des articles~\cite{ipsum1998foundational} ou à des ressources en ligne~\cite{placeholder_www}, munies d'une icône cliquable menant à la page concernée. Comme pour les acronymes, vous pouvez survoler la référence pour en avoir un aperçu.
 ]

\section[demand_research_global]{Vision globale}

Lorem ipsum dolor sit amet, consectetur adipiscing elit, sed do eiusmod tempor incididunt ut labore et dolore magna aliqua.
Mes travaux portent sur l'\acrofull{ai_fr}, et plus particulièrement sur les objets fictifs~\cite{ipsum1998foundational}.

Ut enim ad minim veniam, quis nostrud exercitation ullamco laboris nisi ut aliquip ex ea commodo consequat~\cite{lastname2021journal}.
Duis aute irure dolor in reprehenderit in voluptate velit esse cillum dolore eu fugiat nulla pariatur~\cite{amet2017seminal}.

\begin{boxenv}[demand_research_question]{Question}[Formuler le problème en une phrase]
    Les boîtes fonctionnent de la même manière que dans le manuscrit complet.
    Le type est libre : ici \inquotes{Question}, ailleurs \inquotes{Défi} ou \inquotes{Définition}.
\end{boxenv}

\section[demand_research_results]{Principaux résultats}

Les résultats de \nameref{demand_research_figure} illustrent le comportement attendu d'un modèle fictif.
Le détail des travaux correspondants est disponible dans \cite{lastname2022conference} et \cite{lastname2024journal}.

\begin{figure}[demand_research_figure]{Un exemple de figure dans le document court. La numérotation des figures est continue sur tout le document, ce qui rend les renvois comme \texttt{\textbackslash nameref\{demand\_research\_figure\}} non ambigus.}
    % Example of a vector figure, meant to be included with \input{} inside a figure environment
% Such fragments are typically produced by a script, so that data and rendering stay in sync
% The tikzpicture environment is patched by the template to never overflow \linewidth

\begin{tikzpicture}

    \begin{axis}
        [
            height      = 5cm,
            width       = \linewidth,
            xlabel      = {Number of training epochs},
            ylabel      = {Placeholder metric},
            xmin        = 0, xmax = 20,
            ymin        = 0, ymax = 1,
            grid        = major,
            grid style  = {separatorColor},
            legend pos  = south east,
            legend cell align = left,
        ]

        \addplot[titlesColor, line width=1.5pt, mark=*, mark size=1.5pt] coordinates
            {
                (0, 0.10) (2, 0.38) (4, 0.55) (6, 0.66) (8, 0.73)
                (10, 0.78) (12, 0.82) (14, 0.84) (16, 0.86) (18, 0.87) (20, 0.87)
            };
        \addlegendentry{Placeholder model}

        \addplot[linksColor, line width=1.5pt, dashed, mark=square*, mark size=1.5pt] coordinates
            {
                (0, 0.10) (2, 0.30) (4, 0.44) (6, 0.53) (8, 0.59)
                (10, 0.63) (12, 0.66) (14, 0.68) (16, 0.69) (18, 0.70) (20, 0.70)
            };
        \addlegendentry{Placeholder baseline}

    \end{axis}

\end{tikzpicture}

\end{figure}

\begin{equation}[demand_research_equation]
    \mathcal{L}(\theta) = \frac{1}{N} \sum_{i = 1}^{N} \left\| f_\theta(\bm{x}_i) - \bm{y}_i \right\|_2^2 + \lambda \left\| \theta \right\|_2^2
\end{equation}

L'\nameref{demand_research_equation} donne la fonction de coût minimisée, où $\bm{x}$ est une entrée, $\bm{y}$ sa cible, et $\theta$ les paramètres du modèle.

\section[demand_research_future]{Perspectives}

Excepteur sint occaecat cupidatat non proident, sunt in culpa qui officia deserunt mollit anim id est laborum.
Ces perspectives prolongent \cite{lastname2025preprint}, et s'appuient sur les ressources décrites dans \cite{placeholder_www}.

\begin{itemize}
    \item Une première direction fictive, sur les données fictives.
    \item Une deuxième direction fictive, sur les modèles fictifs.
    \item Une troisième direction fictive, sur les applications fictives.
\end{itemize}

%%%%%%%%%%%%%%%%%%%%%%%%%%%%%%%%%%%%%%%%%%%%%%%%%
%%%%%%%%%%%%%%%%%% BIBLIOGRAPHY %%%%%%%%%%%%%%%%%
%%%%%%%%%%%%%%%%%%%%%%%%%%%%%%%%%%%%%%%%%%%%%%%%%

% In a short document, the list of publications is a section of the current chapter
\printbibliography
